\begin{prob}[source={2022 年全国新高考 I 卷第 18 题}, topic={取值范围}]
记 $\triangle ABC$ 的内角 $A,B,C$ 的对边分别为 $a,b,c$, 已知 $\frac{\cos A}{1+\sin A} = \frac{\sin 2B}{1+\cos 2B}$.
\begin{enumerate}
	\item[(1)] 若 $C = \frac{2\pi}{3}$, 求 $B$；
	\item[(2)] 求 $\frac{a^2+b^2}{c^2}$ 的最小值。
\end{enumerate}
\end{prob}
\begin{prob-sol}
本题的核心数学对象是 $\triangle ABC$，回顾我们所学的知识，一个三角形的\textbf{形状}由其三个内角 $A, B, C$ 决定.由于内角和为 $\pi$，即 $A+B+C=\pi$，，是为其一之约束，因此，一个三角形的形状本质上只有 \textbf{2 个自由度}.我们可以选择任意两个角（例如 $A$ 和 $B$）作为独立的参数，第三个角 $C$ 便随之确定.

题目给出的条件 $\frac{\cos A}{1+\sin A} = \frac{\sin 2B}{1+\cos 2B}$ 是一个满足式条件，而我们的首要任务，就是将这个不咋好看的等式，翻译成一个关于角 $A, B$ 的更简洁的关系式.

因此，我们化简下原式，在等式左边，利用半角公：
\[
\frac{\cos A}{1+\sin A} = \frac{\sin(\frac{\pi}{2}-A)}{1+\cos(\frac{\pi}{2}-A)} = \frac{2\sin(\frac{\pi}{4}-\frac{A}{2})\cos(\frac{\pi}{4}-\frac{A}{2})}{2\cos^2(\frac{\pi}{4}-\frac{A}{2})} = \tan\left(\frac{\pi}{4}-\frac{A}{2}\right)
\]
等式右边，利用二倍角公式：
\[
\frac{\sin 2B}{1+\cos 2B} = \frac{2\sin B \cos B}{1+(2\cos^2 B-1)} = \frac{2\sin B \cos B}{2\cos^2 B} = \tan B
\]
因此：
\[
\tan\left(\frac{\pi}{4}-\frac{A}{2}\right) = \tan B
\]
由于 $A, B \in (0, \pi)$, 可知 $\frac{\pi}{4}-\frac{A}{2} \in (-\frac{\pi}{4}, \frac{\pi}{4})$.在各自的定义域内，正切函数是单调的，故：
\[
\frac{\pi}{4}-\frac{A}{2} = B \quad \implies \quad A+2B = \frac{\pi}{2}
\]
这个简洁的关系，就是原复杂三角等式背后隐藏的约束.它将三角形形状的 2 个自由度削减为了 \textbf{1 个自由度}.接着，我们只需要确定一个角，整个三角形的形状就完全确定了.

\textbf{对于第 (1) 问：}
题目给出了构造式条件 $C = \frac{2\pi}{3}$，这提供了消除最后一个自由度的信息.我们现在有两个关于 $A, B, C$ 的线性约束：
\begin{align*}
	A+B+C &= \pi \quad  \\
	A+2B \qquad &= \frac{\pi}{2} \quad
\end{align*}
将 $C=\frac{2\pi}{3}$ 代入第一个式子，得到 $A+B = \frac{\pi}{3}$.联立方程组：
\[
\begin{cases}
	A+2B = \frac{\pi}{2} \\
	A+B = \frac{\pi}{3}
\end{cases}
\]
解得 $B = \frac{\pi}{6}$.至此，三角形的所有角都被唯一确定，问题解决.

\textbf{对于第 (2) 问：}
此问没有给出额外条件，因此我们需要在由 $A+2B=\frac{\pi}{2}$ 所确定的\textbf{所有可能}的三角形形状中，寻找目标表达式的最小值.这正是用自由度参数表达并求解思想的用武之地.

我们选择角 $B$ 作为描述三角形形状的唯一自由参数.其他角可以用 $B$ 表示：
\begin{align*}
	A &= \frac{\pi}{2}-2B \\
	C &= \pi - (A+B) = \pi - \left(\left(\frac{\pi}{2}-2B\right)+B\right) = \frac{\pi}{2}+B
\end{align*}
为了使 $\triangle ABC$ 成立，所有内角必须为正：
\[
A > 0 \implies \frac{\pi}{2}-2B > 0 \implies B < \frac{\pi}{4}
\]
\[
B > 0 \quad \text{且} \quad C = \frac{\pi}{2}+B > 0 \quad (\text{当 } B>0 \text{ 时恒成立})
\]
综上，自由参数 $B$ 的取值范围是 $(0, \frac{\pi}{4})$.

目标函数是 $\frac{a^2+b^2}{c^2}$.利用正弦定理，将其从边长语言翻译为角度语言：
\[
\frac{a^2+b^2}{c^2} = \frac{\sin^2 A + \sin^2 B}{\sin^2 C}
\]
将所有角用参数 $B$ 代入：
\[
f(B) = \frac{\sin^2(\frac{\pi}{2}-2B) + \sin^2 B}{\sin^2(\frac{\pi}{2}+B)} = \frac{\cos^2(2B) + \sin^2 B}{\cos^2 B}
\]
问题转化为：求函数 $f(B)$ 在区间 $B \in (0, \frac{\pi}{4})$ 上的最小值.

为了简化计算，我们进行变量代换.令 $x = \sin^2 B$.
由于 $B \in (0, \frac{\pi}{4})$, 新变量 $x$ 的取值范围是 $(0, \sin^2(\frac{\pi}{4})) = (0, \frac{1}{2})$.
\[
\cos^2 B = 1-\sin^2 B = 1-x \quad \text{且} \quad \cos(2B) = 1-2\sin^2 B = 1-2x
\]
代入得到关于 $x$ 的函数 $g(x)$:
\[
g(x) = \frac{(1-2x)^2+x}{1-x} = \frac{4x^2-3x+1}{1-x}
\]
为求 $g(x)$ 在区间 $(0, \frac{1}{2})$ 上的最小值，我们对其求导：
\[
g'(x) = \frac{(8x-3)(1-x) - (4x^2-3x+1)(-1)}{(1-x)^2} = \frac{-4x^2+8x-2}{(1-x)^2}
\]
令 $g'(x)=0$, 即 $2x^2-4x+1 = 0$.解得 $x = 1 \pm \frac{\sqrt{2}}{2}$.
考虑到定义域 $x \in (0, \frac{1}{2})$，我们取驻点 $x_0 = 1 - \frac{\sqrt{2}}{2}$.
分析可知，当 $x \in (0, x_0)$ 时，$g'(x)<0$；当 $x \in (x_0, \frac{1}{2})$ 时，$g'(x)>0$.因此 $g(x)$ 在 $x=x_0$ 处取得最小值.
将 $x_0 = 1 - \frac{\sqrt{2}}{2}$ 代入 $g(x)$.为简化计算，先对 $g(x)$ 进行代数变形：
\[
g(x) = \frac{-4x(1-x)+(x-1)+2}{1-x} = -4x-1+\frac{2}{1-x}
\]
最小值为：
\begin{align*}
	g\left(1 - \frac{\sqrt{2}}{2}\right) &= -4\left(1-\frac{\sqrt{2}}{2}\right) - 1 + \frac{2}{1-\left(1-\frac{\sqrt{2}}{2}\right)} \\
	&= -4 + 2\sqrt{2} - 1 + \frac{2}{\frac{\sqrt{2}}{2}} \\
	&= -5 + 2\sqrt{2} + 2\sqrt{2} \\
	&= 4\sqrt{2}-5
\end{align*}\hfill\qedsymbol
\end{prob-sol}
